
Standard practice in large language model deployment treats parameter-efficient fine-tuning (PEFT) and KV cache compression as independent, composable modules: the model is fine-tuned first, then a separately-trained eviction policy is applied. This work examines whether the two independently trained components are aligned in their token priority signals. The Spearman rank correlation between LoRA-modified attention weights and Locret contextual importance scores (CIS) is measured at 0.3662 ± 0.0759 across 100 MNLI validation examples on LLaMA-3.1-8B (seed=42; 100\% of examples below the 0.70 alignment threshold), indicating a substantial mismatch between what the task-adapted model attends to and what the LM-loss-trained eviction policy retains.

To address this mismatch, we propose \textbf{JointLoRA-KV}, a training procedure that optimizes LoRA adapter weights and Locret retaining head weights jointly in a single backward pass via task classification cross-entropy loss. A differentiable soft KV budget mask (sigmoid relaxation, temperature τ=0.1) enables gradient flow through the eviction boundary during training; a straight-through estimator (STE) bridges the training-inference discrepancy when hard top-k eviction is restored at inference time. Separate learning rate groups are used for LoRA (1×10⁻⁴) and Locret (5×10⁻⁴) parameters.

Three experiments are reported. (1) The misalignment between task-adapted LoRA attention and LM-trained Locret CIS is confirmed: mean Spearman ρ = 0.3662, 100/100 examples below threshold (H-E1, PASS). (2) In a proof-of-concept run on LLaMA-3.1-8B (1 seed, 1 epoch, 500 training samples), joint training improves mean GLUE accuracy by +1.50 percentage points over a frozen-Locret baseline (B1: 44.00\% → 45.50\%); the pre-registered gate threshold was ≥2.0pp, so this result is classified PARTIAL (H-M1). Gradient flow to Locret retaining heads was confirmed (locret_grad_received=True, grad norms up to ~2.5×10⁻³). (3) Joint training is stable across 3 random seeds on a small proof-of-concept model (d=64, 2 layers; 0 NaN events, 0 divergence events across seeds 42, 123, 456; H-M2, PASS). The primary performance claim — that JointLoRA-KV improves over the sequential fine-tuning baseline (B3) by ≥3\% on LongBench-QA at 50\% KV budget — was not tested; the corresponding experiment (H-M3) was not executed due to compute constraints.

